\documentclass[UTF8,a4paper,12pt]{ctexart}
\usepackage[margin=2.4cm]{geometry}
\usepackage{setspace}
\usepackage{longtable,booktabs,array}
\usepackage{hyperref}
\usepackage{enumitem}
\hypersetup{colorlinks=true,linkcolor=blue,citecolor=blue,urlcolor=blue}
\setstretch{1.18}
\title{从生成归因到人类主体性证据：人机协作贡献溯源的研究进展与缺口}
\author{}
\date{2026年8月}

\begin{document}
\maketitle

\begin{abstract}
生成式人工智能和编程代理正在把创作活动从“人写、机器执行”转变为持续的混合工作流。现有研究主要沿三条路径推进：用检测器或水印判断内容是否由模型生成，用版本控制和操作日志描述产物如何变化，用治理框架要求披露、监督和问责。这些路径分别提供了内容来源、系统状态或制度责任的局部证据，却很少能够重建人类提出问题、界定目标、选择方案、评价结果和作出最终授权的认知链条。本文综合人机协作、软件与数据溯源、生成内容归因、可审计日志和负责任人工智能研究，区分“产物作者”“操作执行者”和“贡献主体”三个概念，并分析证据粒度、语义解释、隐私保护、跨工具互操作和实证评测方面的不足。综述表明，下一步需要一种以事件为基本单位、同时记录 Human、AI、Tool 和 System 来源并绑定项目上下文的贡献溯源模型；其目标不是给出武断的人类贡献百分比，而是形成可验证、可审计且能解释人类主体性的过程证据。一个可行的研究议程包括开放格式、威胁模型、跨任务数据集、与 Git/检测器/普通日志的对照实验，以及面向开发者和评审者的可理解性研究。
\end{abstract}

\section{引言}
当大型语言模型能够规划任务、修改多文件代码、运行测试并根据反馈循环迭代时，传统的作者性假设出现了结构性变化。软件仓库中的最终文件可能由模型生成，但问题由人提出，约束由人设定，方案在多个候选之间由人选择，失败结果由人判断，最终合并也由人授权。类似的链条同样出现在论文写作、数据分析、产品设计和教育评价中。若只问“这段文字或代码是不是 AI 写的”，就会把一个多主体、跨时间的协作过程压缩成二元标签；若只看最终提交者，又会把审阅、决策和责任隐藏起来。

这一变化带来一个比内容检测更基础的问题：如何为人类在 AI 辅助工作中的主体性提供可复核证据。混合智能研究指出，人和 AI 的优势并非简单相加，而是在目标理解、情境判断、创造性探索和自动化执行之间形成互补；同时，过度依赖自动化可能削弱人的元认知监控和决策能力 \cite{hybridintelligencehuman2025,towardascience2026}。有意义的人类监督也不能退化为“点一下接受”，而应包括验证、引导、替换和对外负责等层次 \cite{designingmeaningfulhuman2026,ensuringhumanoversight2023,letmetake2021}。因此，贡献记录的对象不应只是字节或句子，还应包含目标、理由、评价和授权等事件。

既有技术为该问题提供了重要但不完整的基础。数据和工作流溯源研究建立了记录数据流、计算步骤和实体关系的形式模型 \cite{acorecalculus2012,acorecalculus2013,embeddingdataprovenance2016,sharinginteroperableworkflow2019}；可审计日志和防篡改存储研究讨论了时间戳、哈希链、可信执行环境和分布式账本 \cite{proggeranefficient2014,trustbutverify2023,teepatee2024,usingblockchainledgers2025}；生成内容研究则集中在水印、分类器和作者归因 \cite{authorshipattributionin2025,machinegeneratedtext2023,distinguishinghumanfrom2024,curvemarkdetectingai2025}。这些工作分别回答“数据从哪里来”“系统做了什么”和“内容像不像模型生成”，但尚未形成统一的、跨工具的人机贡献表示。

本文的中心主张是：人机协作研究应从“检测 AI 产物”转向“表示和解释贡献过程”。为此，本文首先界定贡献溯源与相邻概念的边界，随后从人类主体性、技术记录、归因检测、治理安全五个方面总结进展，最后提出面向 BAC 类系统的可检验研究议程。全文把 tamper-evident 理解为“篡改可被发现”，而非不可修改；把记录视为审计辅助材料，而非自动裁判学术署名、法律所有权或责任归属。

\section{概念边界：从作者标签到贡献事件}
\subsection{三个容易混淆的对象}
“作者”至少包含三种不同关系。第一种是产物作者，指文本、代码或图像的可见生成来源；第二种是操作执行者，指实际调用模型、编辑文件或运行命令的主体；第三种是贡献主体，指对问题定义、方案选择、质量判断和最终责任产生实质影响的主体。传统 Git 提交主要记录操作执行者和快照差异，不能可靠地区分手写、模型生成、复制粘贴或代理自主修改。水印和检测器关注产物作者，却通常无法回答人类如何设计提示、筛选候选或重写结果。治理文件则提出应由人负责，但常缺乏能够支持事后复核的过程证据。

贡献溯源需要把这三种关系分开编码。一个“AI 生成代码”事件可以指向一组文件差异和模型版本；随后的人类审阅事件可以记录测试标准、修改意见和批准范围；工具事件则保存命令、退出码、测试摘要与文件哈希；系统事件记录时间、会话和账本状态。这样，接受模型输出不会被错误地重写为人类创作，人工改写也不会被粗略地归入模型贡献。该分离与可解释自动决策研究中从审计轨迹生成面向不同受众解释的思路一致 \cite{addressingregulatoryrequirements2020,decisionprovenanceharnessing2018}。

\subsection{贡献不是百分比}
把人类贡献压缩为一个百分比具有直观吸引力，却缺少稳定的测量单位。十分钟的关键架构判断可能比数小时的代码输入更重要；一次否决错误方案可能改变整个项目路径；提出约束和定义验收标准也难以映射为行数。混合智能和人机组队研究因此强调互补关系、动态控制和共享目标，而不是单一生产率指标 \cite{towardascience2026,hybridintelligencehuman2025,prioritizinghumanai2026}。贡献溯源的目标应是保留可解释的证据集合，让评审者根据场景和制度规则判断贡献，而非制造看似精确但不可辩护的分数。

\subsection{项目上下文是证据的一部分}
同一个事件在不同项目中的含义不同。一个代码补丁可能是从零生成，也可能是对既有人工设计的机械重构；一个测试通过可能证明功能，也可能只覆盖了很小的输入空间。因此，记录需要绑定仓库状态、文件路径、父事件摘要、工具和模型上下文，并允许将事件与需求、问题、决策和验证证据关联。工作流和科学计算领域已经表明，缺少上下文和跨平台互操作会使 provenance 难以复现和解释 \cite{sharinginteroperableworkflow2018,sharinginteroperableworkflow2019,recordingprovenanceof2024}。在人机协作中，项目上下文还必须处理私有提示、第三方材料和敏感凭据的最小披露。

\section{人类主体性与人机协作}
\subsection{监督的实质是评价和控制}
“人在回路”常被当作安全口号，但不同研究对人的位置和权力描述并不一致。教育和医疗研究将解释、审计轨迹、双重评分和专业培训视为责任链的一部分 \cite{humaninthe2025,prioritizinghumanai2026}；可变自主性研究则展示了在不同任务阶段让人接管、暂停或撤销系统的重要性 \cite{letmetake2021}。有意义的监督要求人拥有足够的信息、时间和权限去发现错误，并能将判断转化为系统状态变化。如果界面只展示最终答案而不显示候选、输入、约束和失败尝试，形式上的批准并不能证明主体性。

设计有意义监督的最新框架把 AI 的“操作性主体性”和人的“评价性主体性”分开，强调与外部标准相连的解释，而不是要求用户理解模型内部的不可验证推理 \cite{designingmeaningfulhuman2026}。这为贡献溯源提供了直接启示：账本应记录人评价了什么、依据何种标准、是否改变了方案，以及在何处行使了替换权，而不是记录一段被误认为思维链的私密文本。对于代理式编程，尤其需要区分代理提出的计划、代理执行的命令、工具返回的事实和人类最终授权。

\subsection{协同创作中的认知贡献不可见}
混合智能研究反复指出，人类擅长长期目标、情境理解、价值权衡和伦理判断，而现有模型在因果理解、适应新情境和责任承担方面仍有限 \cite{hybridintelligencehuman2025,towardascience2026}。设计领域的混合主动协作也表明，可解释性和交互控制会改变创作者对 AI 建议的采用方式 \cite{generativeaiin2025}。但现有用户研究多以任务完成时间、产出质量、满意度或信任量表为终点，很少将“提出关键问题”“拒绝不合适的建议”“建立评价标准”作为可分析的事件。

教育和学术诚信研究进一步显示，自动化可能导致认知卸载、过度依赖和能力退化；同时，简单的 AI 使用披露并不能说明人是否真正理解和审查了结果 \cite{artificialintelligenceimplications2023,ethicalproblemsin2025,asurveyon2023}。因此，贡献溯源应支持对人类活动的正向表示，而不是把人类仅当作模型输出的被动批准者。研究需要新的任务范式：让参与者完成同一目标但采用不同的代理权限，比较事件记录能否恢复其策略、判断和责任边界。

\subsection{需要跨模态、跨角色的表示}
真实工作流包含文字提示、语音指令、屏幕操作、代码差异、命令输出、测试报告和口头审阅。若只记录文本提示，会漏掉口述需求和手工编辑；若记录全部原始内容，又可能暴露私有数据和密钥。BAC 类系统可以采用摘要、哈希、输入通道和经过脱敏的证据组合，保存“发生过什么”的可验证线索而不复制完整隐私内容。贡献记录还应允许多个角色共同参与：人类需求方、开发者、评审者、AI 模型、编程代理、版本控制器和测试工具各有不同来源类型。这样的角色分离是避免 source laundering（把模型工作重新声明为人工工作）的前提。

\section{技术溯源与可审计记录}
\subsection{形式模型提供了可组合的基础}
Provenance 领域长期研究实体、活动和生成关系，并发展了可重放、切片和安全属性。核心演算把程序执行轨迹作为可计算对象，讨论了披露与混淆的形式性质 \cite{acorecalculus2012,acorecalculus2013}。学习健康系统的实践表明，自动捕获数据处理轨迹可以支持复现、监管合规和决策审计，但必须为具体领域定义有意义的模型，并以最小侵入方式接入既有工具 \cite{embeddingdataprovenance2016}。科学工作流的 CWLProv、RO-Crate 等实践则把运行、输入、输出和环境封装为可交换的包 \cite{sharinginteroperableworkflow2019,recordingprovenanceof2024}。

这些模型的优点是结构清晰、可组合和可查询，局限是大多把活动视为计算步骤，而非带有意图和判断的社会行动。它们可以回答“哪个程序产生了这个文件”，却不一定回答“为什么选择这个程序”“谁定义了通过标准”以及“谁看到了失败结果”。面向数据集责任和授权的工程实践也显示，单独记录数据许可证、来源和版本仍不足以说明谁对模型失败负责 \cite{thedataprovenance2023,towardsaccountabilityfor2021}。BAC 的研究机会在于扩展活动类型：需求、约束、方案、生成、编辑、命令、测试、审阅、批准和撤销都应成为一等事件，并在事件间建立因果或证据关系。

\subsection{日志完整性不等于贡献真实性}
系统安全研究提供了多种 tamper-evident 机制。Progger 通过内核空间记录提高粒度，并处理伪造条目、跨机器时间同步和日志增长问题 \cite{proggeranefficient2014}；临床 AI 日志用密码时间戳、哈希链和 Merkle 树绑定数据流与推理结果 \cite{trustbutverify2023}；可信执行环境和区块链进一步讨论远程证明、不可抵赖和分布式存储 \cite{teepatee2024,usingblockchainledgers2025,blockchainbased2024}。这些方法能发现记录被修改，却不能保证原始记录者诚实，也不能证明事件描述完全反映了人的真实意图。

因此，安全边界必须诚实表达。哈希链能发现插入、删除、重复和重排，但若攻击者控制写入端，仍可能从一开始就写入虚假事件；区块链能提高共享账本的可验证性，却带来隐私、成本、吞吐和治理问题 \cite{systematicreviewon2022,thesecurityand2015}。贡献溯源需要把完整性证据与来源声明、工具观测和人工确认分开，形成“谁声称了什么、哪个工具观察到了什么、哪些内容仍未验证”的分层结构。

\subsection{互操作性和可查询性仍然不足}
数据溯源标准往往针对科学工作流、数据库或安全取证中的特定语义。跨平台交换、版本兼容、事件粒度和隐私策略缺乏统一约定，导致一个系统的 provenance 很难被另一个系统读取和比较 \cite{sharinginteroperableworkflow2018,traceabilityfortrustworthy2021}。软件开发中的 Git、CI、IDE、代理主机和云执行器各自拥有日志，却没有共同的贡献事件词汇。研究还常将日志设计为事后取证工具，而非在工作时辅助用户理解和控制代理。

一个开放的贡献格式至少应提供：稳定的事件类型；Human、AI、Tool、System 来源枚举；项目、仓库和文件锚点；前序摘要与内容哈希；时间和会话上下文；脱敏证据；验证状态；以及可扩展的签名和外部时间锚。事件应追加而非原地覆盖，并允许针对不同审计者生成最小披露视图。面向可信 AI 的 provenance 文献综述也把可解释性、数据血缘和验证接口视为彼此耦合的设计问题，而不是单一日志字段 \cite{provenancedocumentationto2022,traceabilityfortrustworthy2021}。这样的设计可把形式 provenance、软件工程证据和人机交互记录连接起来。

\section{生成内容归因与检测的边界}
\subsection{检测回答的是来源概率}
AI 生成文本检测、作者归因和水印研究发展迅速。方法包括风格特征、神经分类器、模型概率曲率、人类辅助判断和语义水印 \cite{machinegeneratedtext2023,cognitiveconstraintsimulation2023,distinguishinghumanfrom2024,curvemarkdetectingai2025}。它们可在受控数据上区分某些模型输出与人类样本，也能帮助平台发现批量生成或恶意内容。然而，检测结果通常是对一段成品的概率判断，无法恢复生成前后的人类编辑、提示策略、拒绝的候选和评价标准。

当人和模型交替修改文本时，二元分类尤其脆弱。轻微改写、翻译、混合复制和模型更新会改变统计特征；水印可能在重写、截断或跨平台传输后失效，也可能与隐私和表达自由发生冲突。作者归因综述明确指出，LLM 模糊了人类作者与机器作者的界线，现有任务定义和数据集仍不稳定 \cite{authorshipattributionin2025}。因此，“检测到 AI”不能直接推导“人类没有贡献”，“未检测到 AI”也不能证明完全由人类创作。

\subsection{内容 provenance 仍缺少过程语义}
面向数字出版的基础设施研究提出用不可变存储、隐私保护身份和数据 provenance 支持作者性和责任 \cite{rethinkingdatainfrastructure2022}；AIGC 调查则从模型、内容和安全威胁角度整理水印与溯源技术 \cite{asurveyon2023,asurveyon20231}。这些工作为内容认证提供了重要组件，但通常将 provenance 理解为文件来源或平台签发记录，尚未编码“人类是否提出目标”“人类是否理解并验证结果”等主体性问题。

对 BAC 而言，检测器更适合作为外部交叉证据，而不是账本的核心。一个事件可以引用检测结果、模型声明或平台签名，但必须保留其适用范围、置信度和潜在误报。针对数据取证和入侵调查的 provenance 系统说明，图结构和透明日志可以增强事件关联，但也会产生存储膨胀、查询复杂度和误报权衡 \cite{forensiblockaprovenance2023,provenancebasedintrusion2022,omegaloghighfidelity2020}。更有价值的评测不是只测检测准确率，而是测试审计者能否依据事件链解释贡献边界，并在有意的 source laundering、遗漏日志、重排事件和隐私删减条件下发现风险。

\section{治理、责任与隐私安全}
\subsection{制度要求正在从披露走向可审计}
负责任 AI 研究已经形成透明、公平、隐私、问责和人类监督等原则集合，但原则落地存在工具碎片化和情境不匹配问题 \cite{connectingthedots2023,puttingaiethics2021,aiethicsintegrating2025}。算法问责框架强调组织流程、影响评估和第三方审计；外部监督研究提醒，审计生态需要独立角色、访问权限、证据保存和申诉机制，而非只依赖开发者自报 \cite{agovernanceframework2019,outsideroversightdesigning2022}。监管解释研究表明，从 provenance 生成面向消费者、审计员和监管者的不同解释是可行方向 \cite{addressingregulatoryrequirements2020}。

这些研究为贡献账本提出双重要求。第一，记录必须能支持具体制度问题，例如模型是否被使用、谁批准了部署、哪个测试证明了安全边界。医疗 AI 和 AMIA 原则等高风险场景进一步说明，透明性必须和可追溯、数据准备、临床责任及部署后的持续监测一起设计 \cite{anauditableand2026,definingamias2021}。第二，记录不能把复杂责任链简化为“有日志即合规”。治理还涉及权力、资源、申诉、专业判断和组织文化。BAC 可以成为证据基础设施，但最终裁决仍需由项目政策、机构规则和人类评审完成。

\subsection{隐私最小化与可验证性存在张力}
完整提示、屏幕录制和原始命令可能含有个人信息、商业机密、第三方文本和凭据；过度记录会扩大攻击面，过度删减又使贡献无法复核。信息问责研究说明，使用政策和 provenance 日志本身也需要访问控制、隐私保护和安全建模 \cite{thesecurityand2015}。数据治理和数据信托研究提出让数据主体查看使用路径、控制访问和保留审计权，但跨组织共享仍面临互操作和信任成本 \cite{datatrustsas2022,establishingethicalframeworks2025}。

可行的设计原则包括：默认记录摘要和哈希而非完整私有提示；将密钥、令牌和无关个人数据列为禁止字段；对不同角色提供分级视图；把删除或修订记录为新的披露事件；在需要真实性时使用签名或外部可信锚，而不是只相信声明字段。隐私保护不是把事件链变成黑箱，而是明确哪些事实可验证、哪些事实被有意隐藏、谁有权解密以及隐藏会造成何种证据缺口。

\subsection{攻击面需要专门威胁模型}
现有日志研究关注伪造、删除、重排和拒绝服务；AI 治理研究关注偏差、滥用、过度依赖和责任转移。贡献溯源还面临新的 source laundering：人类把 AI 生成描述为手写，代理把工具失败伪装成成功，平台截断不利的历史，或用户只提交有利于自己的账本视图。攻击者还可能通过短提示哈希猜测内容、利用时间戳回滚或在多个账本间复制事件。

因此，系统评测应将完整性、来源真实性、语义充分性和隐私泄露分别测量。一个哈希链验证通过的账本可能仍缺少关键人工决策；一个内容检测器准确率很高的系统可能无法恢复贡献；一个记录极其详尽的系统可能不可部署。将这四类性质分开，是避免夸大“不可篡改”或“自动判定作者”的必要条件。

\section{讨论：现有研究的共同缺口}
综合以上证据，可以看到五个相互关联的缺口。第一，\textbf{对象缺口}：研究分别记录文件、数据流、模型输出或治理流程，缺少统一的人机贡献事件本体。第二，\textbf{语义缺口}：日志能说明发生了哪些操作，却不能表示目标、理由、评价标准和授权范围。第三，\textbf{主体缺口}：人类常被建模为操作者或批准者，提出问题、构造约束、拒绝建议和承担责任等高价值活动没有稳定表示。第四，\textbf{互操作缺口}：Git、IDE、代理主机、CI 和审计平台的记录难以拼接，跨工具因果链容易断裂。第五，\textbf{证据缺口}：多数工作报告原型或模拟性能，缺少与 Git、普通日志和 AI 检测器的直接对照，也缺少真实协作任务中对“是否正确恢复贡献”的测量。

这些缺口解释了为什么单纯增加日志数量不能解决问题。更细粒度的系统调用记录可能提高取证能力，却不一定让评审者理解人类的决策；更强的水印可能提高内容认证，却无法表示人类如何改写；更严格的披露规则可能提高制度透明度，却可能产生形式化勾选和新的隐私风险。真正的研究挑战是把事件链设计成一种可读的中介：它既能被机器验证，也能让人解释协作如何展开。

BAC 提供了一个具有研究价值的工程化假设：以追加式、项目绑定的 `.bac` 容器记录 Human、AI、Tool、System 四类来源，并用规范化 JSON、哈希链、文件证据、命令结果和验证锚点把事件连起来。该假设仍需通过实证检验，而不能由格式本身推出有效性。尤其需要验证：账本是否能减少贡献争议；不同经验的审计者是否能据此一致恢复关键人类决策；记录负担是否可接受；脱敏后证据是否仍足够；以及攻击者能否利用格式漏洞制造误导性历史。

\section{展望：可检验的研究议程}
\subsection{贡献事件本体与开放格式}
第一步是形成跨任务的最小本体，至少区分需求、约束、计划、AI 生成、工具执行、人工编辑、审阅、测试、批准、撤销和外部锚定。每个事件应有来源类型、时间、上下文、前序摘要、目标对象、证据引用和验证状态。格式应支持 JSON Schema、版本演进和最小披露视图，并与 W3C PROV、软件供应链 attestation、CWLProv 和 Git 对照映射，而不是重新发明不可互操作的术语。

\subsection{基准任务与因果恢复评测}
第二步是建立包含真实或高保真模拟协作过程的数据集。每个任务保留完整的隐藏真值：人类需求、候选方案、模型输出、工具结果、人工修改和最终批准；同时生成 Git-only、普通操作日志、内容检测器和 BAC 等不同证据条件。评测指标不应只有准确率，还应包括关键事件召回、因果顺序恢复、贡献边界解释一致性、审计者完成时间、误归因率、隐私泄露率和记录开销。可设计 source laundering、删尾、重排、代理越权和提示哈希猜测等对抗情境。

\subsection{用户研究与人类主体性}
第三步是开展开发者、评审者和治理人员的混合方法研究。定量实验可以比较有无账本时的责任判断、信任校准和错误发现；访谈和屏幕回放可以考察用户如何理解“提出问题”“选择路径”和“最终授权”等高层事件。研究还应覆盖文字、代码、语音和多模态输入，比较不同记录粒度对认知负担和主体性感受的影响。关键问题不是用户是否喜欢日志，而是他们能否用日志做出更准确、更公平且可解释的判断。

\subsection{隐私、密码与组织治理}
第四步是把技术安全和制度安全同时纳入设计。需要评估哈希链、数字签名、可信时间戳和私有锚点在密钥丢失、时钟漂移、分叉合并和离线工作下的行为；研究差分披露、字段级加密和可撤销访问如何影响审计充分性；并明确谁维护账本、谁承担校验失败责任、谁有权发起争议。只有当部署成本、隐私边界和治理权责被共同定义，贡献账本才可能从演示原型走向可用基础设施。

\subsection{从软件工程扩展到知识生产}
最后，应验证模型是否可迁移到论文、数据分析、设计和教育评价。可负责的数据分析系统已经展示了把文件哈希、版本控制和可复现计算关联起来的可行性，这为跨领域事件映射提供了经验起点 \cite{asystemfor2018}。不同领域对“贡献”的规范不同：软件更重视版本和测试，论文更重视问题、论证和引用，设计更重视迭代与审美判断，教育更重视学习者的理解与反馈。跨领域比较可以检验哪些事件是通用的，哪些必须由领域本体扩展；也能避免把软件开发中的行级差异误当作所有创作活动的通用贡献单位。

\section{结论}
人机协作时代的核心问题不是把最终内容分成“人写”或“AI 写”，而是让关键的人类判断、AI 生成、工具观察和系统状态在时间上可追踪、在语义上可解释、在完整性上可验证。现有 provenance、日志、水印和治理研究已经提供了坚实组件，却各自停留在数据流、产物来源或制度原则的局部层面。它们尚未解决贡献主体性表示、跨工具互操作、隐私最小化和因果恢复评测的组合难题。

将 BAC 发展为 Human--AI Contribution Provenance 框架，具有明确的科学问题和工程抓手：定义贡献事件本体，建立开放 `.bac` 格式，构造威胁模型和基准数据集，并通过对照实验检验账本是否比 Git、普通日志和 AI 检测器更准确地恢复和解释真实协作过程。其价值不在于给人类贴上一个百分比标签，而在于为提出关键问题、作出关键判断和组织 AI 完成工作的主体性提供一组诚实、可审计、可争议也可改进的证据。

\bibliographystyle{gbt7714-nsfc}
\bibliography{human_ai_contribution_provenance_参考文献}
\end{document}
